We propose a Model-Assisted Online Learning (MAOL) framework for sequential volatility forecasting and tail-risk management in government bond markets. The framework wraps any parametric volatility model in three adaptive layers driven by the parameter-free Continuous Coin Betting (COCOB) optimizer, requiring no learning-rate tuning. Layer~1 updates base-model parameters online via the QLIKE gradient. Layer~2 applies a multiplicative bias correction $b_t = \exp(\rho_t)$ learned online on the same loss. Layer~3 recalibrates Value-at-Risk in probability integral transform (PIT) space via online gradient descent on the pinball loss. We instantiate the framework with four base models (RV-GARCH, HAR-RV, Log-GARCH, and Log-HAR) and prove a finite-time regret bound for Layers 1-2 and an explicit finite-time unconditional calibration bound for Layer 3, the latter holding pathwise for any deterministic sequence of PIT values without distributional assumptions. We evaluate the framework on 40 years of 10-year US Treasury futures (1983 to 2024). MAOL-Log-GARCH achieves the lowest average regret against the full-sample HAR-RV oracle ($+0.0027$ QLIKE per step), while MAOL-GARCH significantly outperforms EWMA (Diebold-Mariano statistic $-34.5$, $p < 0.001$) and is competitive with standard rolling and expanding HAR benchmarks. The recalibrated 5\% VaR passes all four standard backtests (Kupiec, Christoffersen, Dynamic Quantile, Acerbi-Sz\'ekely) with a violation rate of 4.82\%. Cross-country validation on German Bund, UK Gilt, and Japanese JGB futures confirms consistent outperformance over EWMA across all markets without any market-specific tuning. A sample-length study re-running the framework on non-overlapping windows of 2, 5, 10, and 15 years shows that performance does not depend on a long warm-up period, and a regime analysis reveals that the multiplier responds within days to sharp volatility spikes (COVID-19) while adapting more gradually to prolonged stress episodes (GFC), with negligible spurious activity during tranquil periods.
